% =============================================================================================
% SECTION 1 -- INTRODUCTION.  Draft 2 -- rebuilt on the six-block structure of Paliwal & Alam
% (arXiv:2511.13202v3), which this manuscript follows throughout.
%
% Seven paragraphs, one job each:
%   1 why the property matters          2 the computational route and its cost
%   3 the mismatch, quantified          4 THE GAP -- prior work named, then one shared criticism
%   5 the scarcity that makes it hard   6 what we did, in paper order
%   7 the contrast claim
%
% Draft 1 buried the gap across paragraphs 3, 6 and 7, opened the differentiation defensively
% ("We are not the first to bridge the two scales"), and ended four times over: a numbered
% contributions list, a strategy paragraph, and a pre-emptive disclaimer. The disclaimer moved to
% the Limitations subsection where it can be quantified; the list became prose.
%
% Every citation here was resolved through Crossref from a real search and carries a row in
% claims_ledger.csv. None was written from memory.
% =============================================================================================

\section{Introduction}
\label{sec:intro}

A thermoelectric generator asks one material to carry charge freely and to obstruct heat at the
same time, and half Heuslers are unusually good at the first half of that bargain and unusually
poor at the second. They are mechanically robust, stable to high temperature, and built from
abundant, non-toxic elements; figures of merit approaching $zT \approx 1.5$ have been demonstrated
in $p$-type \ce{FeNbSb}~\cite{fu2015realizing} and $zT \approx 1.4$ in
\ce{ZrCoBi}~\cite{zhu2018discovery}. What holds the class back is that it conducts heat through the
lattice too well. Because the figure of merit varies inversely with the lattice thermal
conductivity $\kappa_L$, essentially every advance in this family has come from suppressing it.
Identifying compounds whose $\kappa_L$ is intrinsically low is therefore the central screening
problem, and it is the one this paper addresses.

In principle that screen can be run without entering a laboratory. Solving the phonon Boltzmann
transport equation from first principles yields $\kappa_L$ with no experimental input, and the
method has been applied across the half-Heusler space to pick out compounds of unusually low
conductivity~\cite{carrete2014finding}. The obstacle is cost. A converged calculation is expensive
enough that exhaustive enumeration is impractical, and a great deal of effort has consequently gone
into reproducing the calculation more cheaply---through compressed-sensing
descriptors~\cite{liu2020a} and, increasingly, through machine-learned models trained on databases
of computed
conductivities~\cite{luo2023predicting,trans2022lattice,miyazaki2021machine,jaafreh2021lattice,bhattacharjee2022thorough,paliwal2025accelerate}.
Those models are fast, and they are accurate against the quantity they were trained on.

That quantity is not the one a laboratory measures. Models in this line are trained on
\emph{calculated} labels because calculated labels are what exist in quantity, so their reported
accuracies are accuracies against calculation; no claim is made in those papers that the output is
a measurement, but the inference is left to the reader and the reader makes it. It does not hold.
In the corpus assembled for this work, published transport calculations exceed the measured lattice
thermal conductivity of the same compound by a median factor of \num{1.55}. A model reproducing
calculations to \SI{10}{\percent} is therefore roughly \SI{55}{\percent} away from the experiment it
exists to guide. That calculated $\kappa_L$ runs high is not itself news---it is documented compound
by compound throughout the half-Heusler literature, and attacked compound by compound by
identifying the responsible defect chemistry. What has not been available is the size of the offset
across a class, and a correction for it validated against held-out experiment.

Three routes across that gap already exist, and the differences between them matter. Zhu
\emph{et al.} trained a crystal-graph neural network on \num{2668} computed conductivities and
transfer-learned it onto \num{132} experimentally measured values, roughly halving the error on the
measured scale; the computed labels they began from were themselves semi-empirical estimates rather
than transport calculations~\cite{zhu2021charting}. Miller \emph{et al.} fitted the parameters of a
semi-empirical Umklapp model directly to experimental data~\cite{miller2017capturing}. Chen
\emph{et al.} bypassed calculation altogether, training a Gaussian-process model on the measured
conductivities of about one hundred inorganic solids~\cite{chen2019machine}. In none of these is the
object being corrected a full Boltzmann-transport calculation. The distinction is not cosmetic:
semi-empirical and transport-calculated conductivities for the same compound differ by a median
factor of \num{1.38} with a tail to \num{7.2} (Section~\ref{sec:tiering}), so they are different
inputs and require different corrections.

Correcting a transport calculation requires measurements to correct it against, and here the
imbalance is severe. Of the \num{742} half Heuslers in our corpus carrying a thermal conductivity
record of any kind, \num{289} have a full transport calculation and \num{52} have ever been measured
in a laboratory---a ratio of \num{5.6} to one. Training a model directly on measurements is the
obvious alternative and is beginning to appear~\cite{barua2024interpreta}, but \num{52} compounds
are too few to learn a structure--property relation from scratch. They are, as we show, enough to
calibrate one.

This paper therefore takes the opposite route to the usual one. Rather than train a model to
reproduce calculations more cheaply, we take the calculations as given and ask what it costs to move
them onto the experimental scale. We assemble an experiment-anchored corpus---\num{52} measured
compounds drawn from \num{164} publications alongside \num{289} carrying transport calculations,
every value tagged with its provenance and an explicit method tier so that measurements and
calculations are never pooled (Section~\ref{sec:dataset}). We measure the offset between the two
scales and fit a two-parameter transfer function that removes it
(Section~\ref{sec:transfer}). We then validate that function with each compound's entire chemistry
withheld from the calibration as well as from the model: inside a domain of applicability declared
in advance, the corrected predictions place \SI{86.8}{\percent} of \num{38} held-out compounds
within a factor of two of measurement at \SI{33.7}{\percent} median error, against
\SI{50.5}{\percent} median error uncorrected and \SI{44.0}{\percent} for the better of two
compound-blind nulls given the same holdout (Section~\ref{sec:blind}). On that basis we issue five
predictions for unmeasured compounds, each with a measured-coverage interval, and decline eleven
further candidates, naming the condition each fails (Sections~\ref{sec:issued}
and~\ref{sec:refusals}). Five modifications to the method were implemented and rejected on held-out
evidence, including one whose premise the data support but whose benefit does not transfer; we
report those alongside the result they failed to improve (Section~\ref{sec:rejected}).

What separates this from the existing routes is the object being corrected and the form of the
correction. The object is the transport calculation itself rather than a semi-empirical proxy for
it. The correction is closed-form---two parameters a reader can apply by hand and interrogate
physically, not a fitted network---and it carries a domain of applicability that was declared
before the validation set was scored and is tested rather than asserted. The result is not a cheaper
calculation. It is a calculation moved onto the scale on which experiments are reported.
